\begin{statementbox}
	\textbf{\textcolor{CStatement}{条件}}
	\begin{description}[style=nextline, leftmargin=2.8em, labelsep=0.5em, itemsep=0.35em]
		\item[\textbf{\textcolor{CStatement}{条件1（几何可补）：}}] 几何体可以补形为一个长方体，且补得的长方体共顶点三条棱长 $a,b,c$ 可确定。
	\end{description}

	\textbf{\textcolor{CStatement}{结论}}
	\begin{description}[style=nextline, leftmargin=2.8em, labelsep=0.5em, itemsep=0.45em]
		\item[\textbf{\textcolor{CStatement}{结论一（外接球半径）：}}]
			\textbf{\textcolor{CStatement}{直观描述：}}
			外接球直径等于长方体体对角线长。
			\[
				R = \frac{\sqrt{a^{2}+b^{2}+c^{2}}}{2}
			\]

		\item[\textbf{\textcolor{CStatement}{推广结论（常见情形）：}}]
			\textbf{\textcolor{CStatement}{直观描述：}}
			常见可补形几何体包括墙角模型、对棱相等模型、正四面体及底面为矩形的四棱锥。
			这些情形下，$a,b,c$ 分别对应：
			\begin{itemize}
			\item 墙角模型：$a,b,c$ 即为三条两两垂直的棱长；
			\item 对棱相等模型：$a,b,c$ 满足各组对棱分别为长方体相对面的面对角线；
			\item 正四面体：补形得正方体，$a=b=c$；
			\item 线面垂直且底面为矩形的四棱锥：$a,b,c$ 对应矩形的边长和高。
		\end{itemize}
	\end{description}
\end{statementbox}
